
The framing of bidirectional AI alignment --- the concept that ML/NLP and HCI research communities should mutually engage with one another on alignment problems --- rests on an implicit empirical assumption: that these communities are intellectually coupled in practice. Cross-community citation behavior constitutes a structural indicator of epistemic integration. If researchers from ML/NLP and HCI working on alignment genuinely engage with each other's work, their citation networks should reflect symmetric knowledge flow. If citation is asymmetric --- HCI researchers systematically citing ML/NLP but not vice versa --- then bidirectionality may be more aspirational than structural.

Shen et al. [2024] assembled the huashen218 corpus: a curated systematic bibliography of over 400 interdisciplinary alignment papers spanning NeurIPS, ICML, CHI, CSCW, and adjacent venues. The corpus operationalizes the bidirectional alignment community as a named, bounded set of papers. Measuring directed citation behavior within this corpus is the first empirical test of whether the community is citation-integrated.

Attempting this measurement exposes a deeper infrastructural problem. Standard bibliometric approaches --- classifying papers by venue name substring matching --- achieve only 12\% label coverage on this corpus. The cause is the following: the Semantic Scholar Academic Graph (S2AG), the public API for directed citation data, returns full proceedings names (e.g., ''Proceedings of the 36th Annual Conference on Neural Information Processing Systems'') rather than abbreviated venue strings (e.g., ''NeurIPS''). Substring matching against abbreviated strings fails silently on the long-form variants, leaving 88\% of papers unclassified.

This gap motivates the central technical contribution of the present work. Switching the primary classifier from venue name strings to S2AG's \texttt{fieldsOfStudy} tags --- which encode disciplinary assignment at the knowledge-graph level, independent of venue name formatting --- achieves 100\% paper coverage on the same corpus. This is a three-line implementation change with an eight-fold coverage improvement.

The paper makes four contributions:

\begin{enumerate}
\item \textbf{A validated S2AG bibliometric pipeline} for within-corpus directed citation analysis: batch ID resolution via POST \texttt{/paper/batch}, FoS-primary classification (Scheme 3), within-corpus directed graph construction using NetworkX, and gate-controlled evaluation. The pipeline passes 23/23 unit tests and is fully cached for zero-API re-execution.
\end{enumerate}

\begin{enumerate}
\item \textbf{FoS-primary classification (Scheme 3)}: a classification method using S2AG \texttt{fieldsOfStudy[0]} as the primary venue-group classifier, achieving 100\% label coverage on the huashen218 corpus versus 12\% for venue-string-only approaches. The method is domain-general and applicable to any interdisciplinary corpus indexed by S2AG.
\end{enumerate}

\begin{enumerate}
\item \textbf{The first directional citation measurement within the huashen218 alignment corpus}: in the 33-paper resolved proxy, ML/$NLP\rightarrow HCI$ proportion = 10.6\%, $HCI\rightarrow ML/NLP$ proportion = 100\%, ratio $\approx$ 0.107. No statistical falsifier was triggered; full statistical testing requires an augmented corpus (designated h-e1-v2 in the research log).
\end{enumerate}

\begin{enumerate}
\item \textbf{A methodological finding on corpus proxy construction}: the huashen218 GitHub reading list yields 49 extractable identifiers from approximately 130 linked papers --- roughly 12.5\% of the target corpus. Public reading lists from systematic reviews are insufficient bibliometric dataset proxies; programmatic S2AG citation-of-citation search is the viable reconstruction approach.
\end{enumerate}

The paper is organized as follows. Section 2 situates the work within prior cross-community citation analysis and bibliometric infrastructure. Section 3 describes the pipeline design and classification schemes. Section 4 presents the experimental setup. Section 5 reports results. Section 6 discusses implications and limitations. Section 7 concludes.

